% ============================================================
%  APPENDIX A --- Installation
% ============================================================
\chapter{Installation}
\label{app:installation}

La présente annexe couvre tous les scénarios de déploiement soutenus pour FairWindSK,
d'une première source construire sur un ordinateur portable développeur à un entièrement automatique
démarrage de kiosque sur un écran de barre dédié.
Chaque scénario énumère les besoins matériels ou de plate-forme, les commandes exactes,
et les étapes de vérification après l'installation.

\begin{notebox}
Tous les scénarios supposent un serveur Signal~K en marche quelque part sur le bateau
Réseau Wi-Fi.
FairWindSK est un shell d'affichage, pas un serveur de données.
Installez et configurez Signal~K avant d'installer FairWindSK.
\end{notebox}

% ── A.1  Prerequisites ───────────────────────────────────────
\section{Préalables}
\label{app:prerequisites}

\subsection{Prérequis communs pour toutes les plateformes de bureau}

\begin{itemize}
  \item CMake 3.16 or newer
  \item A C++17-capable compiler supported by the chosen Qt~6 kit
  \item Git
  \item Internet access during the first build
        (CMake downloads \code{nlohmann/json}, \code{QtZeroConf}, and \code{QHotkey}
automatiquement; les constructions suivantes utilisent les copies mises en cache)
\end{itemize}

\subsection{Requis Qt 6 modules --- bureau}

\rowcolors{2}{LtGray}{white}
\begin{longtable}{L{4.5cm}L{9.2cm}}
\toprule
\rowcolor{Navy}
\color{white}\sffamily\bfseries Module &
\color{white}\sffamily\bfseries Purpose \\
\midrule
\code{Core}            & Qt base classes and event loop. \\
\code{Gui}             & 2D graphics and font handling. \\
\code{Widgets}         & Widget toolkit used throughout the shell. \\
\code{Concurrent}      & Thread pool for background tasks. \\
\code{Network}         & HTTP/REST client and TLS. \\
\code{WebSockets}      & Signal~K WebSocket stream. \\
\code{Xml}             & Configuration and resource parsing. \\
\code{Svg}             & OpenBridge icon rendering. \\
\code{Positioning}     & GPS location APIs (desktop fallback). \\
\code{WebEngineWidgets}& Embedded browser for hosted Signal~K web apps. \\
\code{VirtualKeyboard} & On-screen keyboard for touchscreen deployments. \\
\code{PrintSupport}    & PDF generation (crash reports). \\
\bottomrule
\end{longtable}

\subsection{Requis Qt 6 modules --- Android et iOS}

Mobile builds replace \code{WebEngineWidgets} with a lighter web layer and
baissez les dépendances de bureau seulement :

\rowcolors{2}{LtGray}{white}
\begin{longtable}{L{4.5cm}L{9.2cm}}
\toprule
\rowcolor{Navy}
\color{white}\sffamily\bfseries Module &
\color{white}\sffamily\bfseries Replaces or adds \\
\midrule
\code{Qml}         & QML runtime for WebView hosting. \\
\code{Quick}       & Qt Quick scene graph. \\
\code{QuickWidgets}& \code{QQuickWidget} bridge between widget and QML worlds. \\
\code{WebView}     & Replaces \code{WebEngineWidgets}. Uses the platform's native browser engine. \\
\bottomrule
\end{longtable}

Mobile builds do \emph{not} include \code{QtZeroConf}, \code{QHotkey},
\code{PrintSupport}, or \code{WebEngineWidgets}.

% ── A.2  Generic desktop build ───────────────────────────────
\section{Construction de bureau générique (toutes les plateformes)}
\label{app:generic}

C'est la procédure de base qui s'applique à chaque plate-forme de bureau
avant d'ajouter des étapes spécifiques à la plateforme.

\begin{enumerate}
  \item Clone the repository:
\begin{lstlisting}
git clone https://github.com/OpenFairWind/FairWindSK.git
cd FairWindSK
\end{lstlisting}

  \item Configure:
\begin{lstlisting}
cmake -S . -B build
\end{lstlisting}

  \item Build (parallel across all CPU cores):
\begin{lstlisting}
cmake --build build --parallel
\end{lstlisting}

  \item Install to the default prefix:
\begin{lstlisting}
cmake --install build
\end{lstlisting}

  \item Override the install prefix if needed:
\begin{lstlisting}
cmake -S . -B build -DCMAKE_INSTALL_PREFIX=/desired/prefix
cmake --build build --parallel
cmake --install build
\end{lstlisting}

  \item If Qt is installed outside the default search path, pass
        \code{-DCMAKE\_PREFIX\_PATH} during configure:
\begin{lstlisting}
cmake -S . -B build -DCMAKE_PREFIX_PATH=/path/to/Qt/6.x.x/gcc_64
\end{lstlisting}
\end{enumerate}

\begin{tipbox}
Si l'accès à Internet est restreint pendant la construction, installer
\code{nlohmann\_json} from your package manager first.
CMake utilisera l'en-tête du système au lieu de le télécharger.
On Debian/Ubuntu: \code{sudo apt install nlohmann-json3-dev}.
On macOS with Homebrew: \code{brew install nlohmann-json}.
\end{tipbox}

% ── A.3  macOS ───────────────────────────────────────────────
\section{installation macOS}
\label{app:macos}

\subsection{Exigences}

\begin{itemize}
  \item macOS 12 Monterey or newer (recommended)
  \item Xcode command line tools: \code{xcode-select --install}
  \item Qt~6 desktop kit (via Qt Online Installer or Homebrew)
  \item CMake and Ninja (via Homebrew or Qt Maintenance Tool)
\end{itemize}

\subsection{Installer avec Homebrew}

\begin{lstlisting}
brew install qt cmake ninja nlohmann-json
cmake -S . -B build -G Ninja \
  -DCMAKE_PREFIX_PATH="$(brew --prefix qt)"
cmake --build build --parallel
\end{lstlisting}

\subsection{Courez à partir de l'arbre de construction}

\begin{lstlisting}
open build/FairWindSK.app
\end{lstlisting}

\subsection{Installation sur le système}

\begin{lstlisting}
cmake --install build
# Installs to /usr/local by default:
open /usr/local/FairWindSK.app
\end{lstlisting}

\subsection{Redistribution (déploiement Qt durée d'exécution)}

Pour la distribution en dehors de l'arborescence de construction, utilisez l'outil de déploiement d'Apple :

\begin{lstlisting}
macdeployqt build/FairWindSK.app
\end{lstlisting}

This bundles all Qt frameworks inside the \code{.app} so the app runs
sur les machines sans Qt installées.

\subsection{Vérification après l'installation}

\begin{enumerate}
  \item Launch FairWindSK.
  \item Open \ui{Settings~> Connection}.
        Enter your Signal~K server address and tap \key{Connect}.
  \item Confirm the status changes to \textit{Live} and apps appear on the launcher.
\end{enumerate}

% ── A.4  Linux desktop ───────────────────────────────────────
\section{Installation du bureau Linux}
\label{app:linux}

\subsection{Exigences (Debian / Ubuntu)}

\begin{lstlisting}
sudo apt update
sudo apt install \
  build-essential cmake ninja-build git \
  nlohmann-json3-dev \
  qt6-base-dev qt6-base-dev-tools \
  qt6-webengine-dev qt6-webengine-dev-tools \
  qt6-websockets-dev qt6-positioning-dev \
  libqt6svg6-dev qt6-virtualkeyboard-dev \
  qt6-base-private-dev libqt6printsupport6 \
  libavahi-compat-libdnssd-dev libnss-mdns avahi-utils
\end{lstlisting}

\subsection{Construire et exécuter}

\begin{lstlisting}
cmake -S . -B build -G Ninja
cmake --build build --parallel
./build/FairWindSK
\end{lstlisting}

\subsection{Installation du système}

\begin{lstlisting}
sudo cmake --install build
\end{lstlisting}

\code{cmake --install} on a Linux desktop target:

\begin{itemize}
  \item Installs the \code{FairWindSK} binary to
        \code{\$\{CMAKE\_INSTALL\_BINDIR\}} (default \code{/usr/local/bin}).
  \item Installs the bundled icon directory alongside the binary.
  \item Installs the desktop-only helper libraries to
        \code{\$\{CMAKE\_INSTALL\_LIBDIR\}}.
  \item Installs the \code{fairwindsk.desktop} launcher file to the standard
Répertoire des applications XDG.
  \item Installs the \code{fairwindsk.png} menu icon to
        \code{/usr/local/share/icons/hicolor/256x256/apps/}.
\end{itemize}

Après l'installation, FairWindSK apparaît dans le menu des applications.
Exécutez-le à partir du menu ou avec :

\begin{lstlisting}
FairWindSK
\end{lstlisting}

\subsection{Emballage}

Inclure le temps d'exécution Qt et les dépendances externes construites sous
\code{build/external/lib} in your package.
The build-tree rpath means the binary runs from \code{build/} without
\code{LD\_LIBRARY\_PATH}, but a packaged install needs those libraries in
le chemin de la bibliothèque système ou groupé à côté du binaire.

% ── A.5  Raspberry Pi OS --- nav station ───────────────────────
\section{Raspberry Pi OS --- Station de navigation}
\label{app:rpi-station}

Ce scénario installe FairWindSK sur un Raspberry~Pi utilisé comme une normale
station de navigation de bureau --- écran, clavier et souris attachés,
exécuter le bureau Raspberry~Pi~OS.

\subsection{Exigences}

Pi~4 ou ~5 de framboise avec Raspberry~Pi~OS Librairie (64-bit recommandé),
et une connexion Internet stable pour le téléchargement initial du paquet.

\subsection{Installer les dépendances de construction}

\begin{lstlisting}
sudo apt update
sudo apt install \
  build-essential cmake ninja-build git \
  nlohmann-json3-dev \
  qml6-module-qt-labs-folderlistmodel \
  qml6-module-qtquick-window \
  qml6-module-qtquick-layouts \
  qml6-module-qtqml-workerscript \
  libnss-mdns avahi-utils \
  libavahi-compat-libdnssd-dev libxkbcommon-dev \
  qt6-base-dev qt6-base-dev-tools \
  qt6-webengine-dev qt6-webengine-dev-tools \
  qt6-websockets-dev qt6-positioning-dev \
  libqt6svg6-dev \
  qt6-virtualkeyboard-dev \
  libqt6virtualkeyboard6 \
  qt6-virtualkeyboard-plugin
\end{lstlisting}

\begin{notebox}
Si la construction échoue avec un manquant
\code{Qt6VirtualKeyboardConfigVersionImpl.cmake}, apply the known
la solution de rechange:
\begin{lstlisting}
sudo touch /usr/lib/aarch64-linux-gnu/cmake/Qt6VirtualKeyboard/\
Qt6VirtualKeyboardConfigVersionImpl.cmake
\end{lstlisting}
Then re-run \code{cmake}.
\end{notebox}

\subsection{Construire et installer}

\begin{lstlisting}
git clone https://github.com/OpenFairWind/FairWindSK.git
cd FairWindSK
cmake -S . -B build -G Ninja
cmake --build build --parallel
sudo cmake --install build
\end{lstlisting}

FairWindSK apparaît dans le menu d'application sous
\textit{Navigation} or \textit{Internet}.

\subsection{Intégration OpenPlotter}

If OpenPlotter is detected during \code{cmake --install}, FairWindSK
copie également son icône de lancement dans les répertoires de menu OpenPlotter.
Aucune étape supplémentaire n'est nécessaire.

% ── A.6  Raspberry Pi OS --- kiosk helm display ────────────────
\section{Raspberry Pi OS --- Kiosk Helm Display}
\label{app:rpi-kiosk}

Voici le scénario recommandé pour un écran de barre dédié:
une framboise ~Pi~4 ou ~5 avec un écran tactile, en utilisant FairWindSK automatiquement
au démarrage sans clavier et sans menu de bureau visible.

\subsection{Aperçu général}

\begin{enumerate}
  \item Install FairWindSK as in Section~\ref{app:rpi-station}.
  \item Configure FairWindSK for full-screen kiosk mode.
  \item Install the \code{fairwindsk-startup} helper to manage autostart.
  \item (Optional) Configure the desktop panel to auto-hide.
  \item Reboot and verify.
\end{enumerate}

\subsection{Étape 1 --- Construire et installer}

Follow Section~\ref{app:rpi-station}.
On Raspberry~Pi~OS, \code{cmake --install} automatically installs the
\code{fairwindsk-startup.desktop} autostart entry to
\code{/etc/xdg/autostart/}.
FairWindSK démarre automatiquement sur la prochaine connexion.

\subsection{Étape 2 --- Configurer le mode kiosque plein écran}

Open \ui{Settings~> Main} and set \key{Window Mode} to
\textbf{Full Screen}.
Cela demande une fenêtre sans cadre, toujours sur le dessus qui couvre l'ensemble
écran sans barre de titre ni décoration de fenêtre.

Alternatively, edit \code{fairwindsk.json} directly:

\begin{lstlisting}
{
  "main": {
    "windowMode": "fullscreen"
  }
}
\end{lstlisting}

\subsection{Étape 3 --- Installez l'aide au démarrage}

The \code{fairwindsk-startup} Python script manages the FairWindSK
le cycle de vie du processus.
It is installed by \code{cmake --install} alongside the binary.

What \code{fairwindsk-startup} does at each boot:

\begin{enumerate}
  \item Reads \code{fairwindsk.ini} to find the configuration path.
  \item Reads \code{fairwindsk.json} to check the \code{virtualKeyboard}
        flag and sets \code{QT\_IM\_MODULE=qtvirtualkeyboard} accordingly.
  \item Waits up to 45 seconds for the Signal~K web-app catalogue
        endpoint to respond (\code{/signalk/v1/apps/list} or the legacy
        \code{/skServer/webapps}).
Cela garantit que les icônes de plugin et les œuvres d'art sont disponibles avant
FairWindSK commence.
  \item Launches FairWindSK.
  \item If FairWindSK exits with code~1 (the \textit{Restart} button in
        \ui{Settings~> System}), the helper relaunches it immediately.
  \item If FairWindSK exits with code~0 (the \textit{Quit} button or
une fermeture propre), l'aide sort proprement.
\end{enumerate}

\begin{warningbox}
Use \key{Restart FairWindSK} in \ui{Settings~> System} (exit code~1)
pour relancer après modification de configuration.
\key{Quit FairWindSK} (exit code~0) leaves the display blank
sans redémarrage.
\end{warningbox}

\subsection{Étape 4 --- Cacher le panneau de bureau (labwc / Bookworm)}

Sur Raspberry~Pi~OS Bookworm avec le compositeur labwc, la barre des tâches
panneau peut superposer la fenêtre plein écran FairWindSK.
Pour éviter cela, configurez le panneau pour cacher automatiquement :

\begin{lstlisting}
# Edit the panel configuration file
nano ~/.config/wf-panel-pi.ini
\end{lstlisting}

Ajouter ou modifier :

\begin{lstlisting}
[panel]
autohide = true
\end{lstlisting}

Enregistrer le fichier et redémarrer.

\subsection{Étape 5 --- Activer le clavier virtuel (écran tactile seulement)}

Si l'écran de barre n'a pas de clavier physique, activez le virtuel Qt
clavier :

\begin{enumerate}
  \item Open \ui{Settings~> Main}.
  \item Tick \key{Virtual Keyboard}.
  \item Tap \key{Restart FairWindSK} in \ui{Settings~> System}.
\end{enumerate}

L'aide au démarrage lit le réglage à chaque démarrage et définit le
variable d'environnement automatiquement.

\subsection{Étape 6 --- Vérifier le démarrage du kiosque}

\begin{enumerate}
  \item Reboot the Raspberry~Pi.
  \item Verify FairWindSK starts automatically (allow up to 60 seconds
pour Signal~K et FairWindSK commencer ensemble).
  \item Confirm the display is full-screen with no title bar.
  \item Confirm the status shows \textit{Live}.
  \item Confirm apps appear on the launcher.
\end{enumerate}

\begin{tipbox}
Sur une framboise ~Pi~4 avec une classe ~A2 carte microSD, FairWindSK
typiquement devient interactif en moins de 30 secondes à partir de la mise en marche,
en supposant que Signal~K commence rapidement.
Use a fast card (\textgreater{}90~MB/s read) to minimise boot time.
\end{tipbox}

% ── A.7  Windows ─────────────────────────────────────────────
\section{Installation de Windows}
\label{app:windows}

\subsection{Exigences}

\begin{itemize}
  \item Windows~10 or~11 (64-bit)
  \item Qt~6 desktop kit for MSVC
(installé via l'installateur en ligne Qt)
  \item Visual Studio Build Tools 2022 or Visual Studio Community 2022
  \item CMake and Ninja
(groupé avec Qt ou installé séparément)
\end{itemize}

\subsection{Construire et installer}

Ouvrez un Developer Command Prompt avec l'environnement MSVC chargé,
puis :

\begin{lstlisting}
git clone https://github.com/OpenFairWind/FairWindSK.git
cd FairWindSK
cmake -S . -B build -G Ninja ^
  -DCMAKE_PREFIX_PATH="C:\Qt\6.x.x\msvc2022_64"
cmake --build build --parallel
cmake --install build
\end{lstlisting}

\subsection{Déployer l'exécution Qt}

Après une compilation réussie, déployez les DLL d'exécution Qt à côté du
exécutable:

\begin{lstlisting}
windeployqt build\FairWindSK.exe
\end{lstlisting}

\code{cmake --install} installs the executable, the bundled icon
directory under \code{bin\textbackslash icons}, and the desktop-only
helper DLLs into the same \code{bin} directory.

\subsection{Cours}

\begin{lstlisting}
build\FairWindSK.exe
\end{lstlisting}

\begin{notebox}
\code{file://} launcher entries are retained in configuration for
compatibilité mais sont bloqués par le modèle à guichet unique sur Windows.
Use \code{http://} or \code{https://} URLs for all applications.
\end{notebox}

% ── A.8  Android ─────────────────────────────────────────────
\section{Installation Android}
\label{app:android}

Android construit utiliser l'implémentation mobile Qt WebView.
Qt Creator est le workflow recommandé pour Android car il gère
la signature SDK, NDK, APK et le déploiement de l'appareil.

\subsection{Exigences}

\begin{itemize}
  \item A macOS, Linux, or Windows build host
  \item Qt~6 with the Android kit for your target ABI
(arm64-v8a est recommandé pour les appareils modernes)
  \item Android Studio (for the SDK, NDK, and emulator)
  \item Java JDK (bundled with Android Studio or installed separately)
  \item CMake and Ninja
\end{itemize}

\subsection{Guide de construction étape par étape}

\begin{enumerate}
  \item Open Android Studio and install:
une plate-forme SDK Android, SDK Platform-Tools, SDK Build-Tools,
et la version NDK recommandée par votre version Qt.
  \item In the Qt Maintenance Tool, install the Android Qt components
qui correspondent à votre hôte et cible ABI.
  \item In Qt Creator, go to
        \ui{Preferences~> Devices~> Android} and verify that the SDK,
Les chemins NDK, JDK et plate-forme sont tous correctement détectés.
  \item Clone the repository:
\begin{lstlisting}
git clone https://github.com/OpenFairWind/FairWindSK.git
\end{lstlisting}
  \item Open \code{FairWindSK/CMakeLists.txt} in Qt Creator.
  \item Select an Android kit (e.g.\ \textit{Android Qt~6.x.x Clang arm64-v8a}).
  \item Let Qt Creator configure the project.
        The mobile configuration uses \code{Qt::WebView},
        \code{Qt::Quick}, and \code{Qt::QuickWidgets},
et saute les dépendances de bureau seulement.
  \item Build with \ui{Build~> Build Project}.
  \item Connect an Android device with developer mode enabled,
ou démarrer un émulateur Android.
  \item Deploy with \ui{Build~> Run}.
\end{enumerate}

\subsection{Configuration en ligne de commande (avancée)}

\begin{lstlisting}
cmake -S . -B build-android -G Ninja \
  -DCMAKE_TOOLCHAIN_FILE="$ANDROID_NDK/build/cmake/android.toolchain.cmake" \
  -DANDROID_ABI=arm64-v8a \
  -DANDROID_PLATFORM=android-24 \
  -DCMAKE_PREFIX_PATH="/path/to/Qt/6.x.x/android_arm64_v8a"
cmake --build build-android --parallel
\end{lstlisting}

Qt Creator est toujours recommandé pour l'emballage et le déploiement de l'appareil.

\subsection{Différences des fonctionnalités Android}

\rowcolors{2}{LtGray}{white}
\begin{longtable}{L{5.2cm}L{8.5cm}}
\toprule
\rowcolor{Navy}
\color{white}\sffamily\bfseries Feature &
\color{white}\sffamily\bfseries Android status \\
\midrule
Hosted web apps          & Supported via Qt WebView. \\
Signal~K stream          & Supported (WebSockets). \\
mDNS server discovery    & Not available (disabled on Android). \\
Global hotkey SHIFT+TAB  & Not available (desktop-only, uses QHotkey). \\
Virtual keyboard         & Uses Android system keyboard; Qt Virtual Keyboard setting has no effect. \\
File printing / PDF      & Not available (PrintSupport not included). \\
QWebEngine cookies       & Not used; mobile path uses a separate mechanism. \\
\code{file://} apps      & Blocked by the single-window model on all targets. \\
\bottomrule
\end{longtable}

% ── A.9  iOS and iPadOS ──────────────────────────────────────
\section{Installation iOS et iPadOS}
\label{app:ios}

iOS et iPadOS construit utilisent la même implémentation web mobile que Android.
Qt Creator avec un compte développeur Apple est nécessaire pour le déploiement de l'appareil.

\subsection{Exigences}

\begin{itemize}
  \item A macOS build host
  \item Xcode and Xcode command line tools
  \item Qt~6 with the iOS kit
  \item CMake and Ninja
  \item An Apple developer account and provisioning profile
(pour le déploiement de l'appareil; pas nécessaire pour le simulateur)
\end{itemize}

\subsection{Guide de construction étape par étape}

\begin{enumerate}
  \item Open Xcode at least once and accept the license agreement.
  \item In the Qt Maintenance Tool, install the iOS Qt components
qui correspondent à votre bureau Qt libération.
  \item Clone the repository:
\begin{lstlisting}
git clone https://github.com/OpenFairWind/FairWindSK.git
\end{lstlisting}
  \item Open \code{FairWindSK/CMakeLists.txt} in Qt Creator.
  \item Select an iOS kit (e.g.\ \textit{Qt~6.x.x for iOS}).
  \item Build with \ui{Build~> Build Project}.
  \item To run on the Simulator: choose an iOS Simulator run target
        and select \ui{Build~> Run}.
  \item To run on a device: connect the iPhone or iPad, select the
périphérique physique dans le sélecteur de cible d'exécution, vérifier les paramètres de signature
quand vous êtes invité, et courir de Qt Creator.
\end{enumerate}

\subsection{Configuration en ligne de commande (avancée)}

\begin{lstlisting}
cmake -S . -B build-ios -G Ninja \
  -DCMAKE_SYSTEM_NAME=iOS \
  -DCMAKE_OSX_SYSROOT=iphoneos \
  -DCMAKE_OSX_ARCHITECTURES=arm64 \
  -DCMAKE_PREFIX_PATH="/path/to/Qt/6.x.x/ios"
cmake --build build-ios --parallel
\end{lstlisting}

Use \code{iphonesimulator} and \code{x86\_64} (or \code{arm64} on
Apple Silicon hosts) pour les constructions Simulator.

\subsection{Différences de fonctionnalités iOS}

\rowcolors{2}{LtGray}{white}
\begin{longtable}{L{5.2cm}L{8.5cm}}
\toprule
\rowcolor{Navy}
\color{white}\sffamily\bfseries Feature &
\color{white}\sffamily\bfseries iOS / iPadOS status \\
\midrule
Hosted web apps          & Supported via Qt WebView (WKWebView). \\
Signal~K stream          & Supported (WebSockets). \\
mDNS server discovery    & Not available. Enter the server URL manually. \\
Global hotkey SHIFT+TAB  & Not available. \\
Virtual keyboard         & Uses iOS system keyboard. \\
File printing / PDF      & Not available. \\
\code{file://} apps      & Blocked by the single-window model on all targets. \\
\bottomrule
\end{longtable}

% ── A.10  Signal K Server on Docker ──────────────────────────
\section{Serveur Signal K sur Docker}
\label{app:signalk-docker}

Si vous n'avez pas déjà un serveur Signal~K, Docker est le plus rapide
C'est un moyen de courir.
Cette section documente le déploiement de Docker
Dépôt de projet FairWindSK.

\subsection{Image Docker}

\begin{lstlisting}
cr.signalk.io/signalk/signalk-server:latest
\end{lstlisting}

Supported architectures: \code{linux/amd64}, \code{linux/arm/v7},
\code{linux/arm64} (covers all Raspberry~Pi~4 and~5 configurations).

\subsection{Démarrer rapidement avec les données de démonstration NMEA}

\begin{lstlisting}
mkdir $HOME/.signalk
docker run --init -it --rm \
  --name signalk-server \
  --publish 3000:3000 \
  -v $HOME/.signalk:/home/node/.signalk \
  cr.signalk.io/signalk/signalk-server \
  --sample-nmea0183-data
\end{lstlisting}

Ceci exécute Signal~K sur port~3000 avec des données NMEA~0183 simulées.
Utilisez-le pour tester FairWindSK sans instruments réels.

\subsection{Déploiement de production}

\begin{lstlisting}
docker run -d --init \
  --name signalk-server \
  -p 3000:3000 \
  -v $HOME/.signalk:/home/node/.signalk \
  cr.signalk.io/signalk/signalk-server
\end{lstlisting}

\begin{itemize}
  \item The \code{-d} flag runs the container as a background daemon.
  \item The \code{-v} flag mounts the configuration directory so
Les paramètres de signal ~K, les plugins et les données persistent à travers les conteneurs
redémarre.
  \item On first access, the Signal~K admin web UI at
        \code{http://localhost:3000} prompts you to create admin credentials.
\end{itemize}

\subsection{Mise à jour du signal K}

\begin{lstlisting}
docker pull cr.signalk.io/signalk/signalk-server:latest
docker stop signalk-server
docker rm signalk-server
# Re-run the production docker run command above
\end{lstlisting}

\subsection{Structure du répertoire à l'intérieur du conteneur}

\rowcolors{2}{LtGray}{white}
\begin{longtable}{L{4.8cm}L{8.9cm}}
\toprule
\rowcolor{Navy}
\color{white}\sffamily\bfseries Path &
\color{white}\sffamily\bfseries Contents \\
\midrule
\code{/home/node/signalk}  & Signal~K server files (read-only inside container). \\
\code{/home/node/.signalk} & Settings files, plugins, and data.
                              \textbf{Mount this to the host.} \\
\bottomrule
\end{longtable}

% ── A.11  First-Run Configuration ────────────────────────────
\section{Configuration de la première course}
\label{app:firstrun}

Après l'installation sur n'importe quelle plateforme, suivez ces étapes lors du premier lancement.

\subsection{Emplacements des fichiers de configuration}

FairWindSK stocke sa configuration dans deux fichiers par utilisateur
répertoire de configuration Qt :

\rowcolors{2}{LtGray}{white}
\begin{longtable}{L{3.6cm}L{10.1cm}}
\toprule
\rowcolor{Navy}
\color{white}\sffamily\bfseries File &
\color{white}\sffamily\bfseries Contents \\
\midrule
\code{fairwindsk.json}  &
Tous les paramètres visibles par l'utilisateur : URL de connexion, mises en page de barres, catalogue d'applications,
widget définitions, préréglages de confort, unités, et cartographies de chemin Signal~K.
JSON lisible par l'homme.
  Safe to back up, transfer, and version-control. \\
\code{fairwindsk.ini}   &
  Qt-managed settings: path to \code{fairwindsk.json}, the debug flag,
et le jeton d'authentification.
Le jeton est gardé ici intentionnellement de sorte qu'il n'est jamais inclus dans un
  configuration export. \\
\bottomrule
\end{longtable}

On first run, if \code{fairwindsk.json} does not exist, FairWindSK
sele it from the packed par défaut in
\code{resources/json/configuration.json}.

\subsection{Se connecter au signal K}

\begin{enumerate}
  \item Open \ui{Settings~> Connection}.
  \item Enter the Signal~K server address,
        e.g.\ \code{http://192.168.1.100:3000}.
Si vous utilisez Signal~K sur le même périphérique,
        use \code{http://localhost:3000}.
  \item Tap \key{Connect}.
  \item Wait for the status to change to \textit{Live}.
  \item If the server requires authentication,
        tap \key{Request Token} and approve the request in the
Interface d'administration Signal~K.
\end{enumerate}

\subsection{Authentifier (facultatif)}

Si le serveur Signal~K a activé la sécurité :

\begin{enumerate}
  \item Tap \key{Request Token} in \ui{Settings~> Connection}.
  \item FairWindSK opens the Signal~K admin interface in the
Tiroir en bas de la barre, navigué vers
        \ui{Security~> Access Requests}.
  \item Approve the request as an admin user.
  \item FairWindSK stores the token automatically and reconnects.
\end{enumerate}

\subsection{Paramètres essentiels pour la première fois}

\rowcolors{2}{LtGray}{white}
\begin{longtable}{L{4.0cm}L{9.7cm}}
\toprule
\rowcolor{Navy}
\color{white}\sffamily\bfseries Setting &
\color{white}\sffamily\bfseries Where and what to set \\
\midrule
Unités&
  \ui{Settings~> Units}.
  Choose knots (speed), nautical miles (distance), and metres or feet (depth). \\
Format des coordonnées&
  \ui{Settings~> Main}.
  Choose DD°MM.MMM' for standard nautical use. \\
Mode fenêtre&
  \ui{Settings~> Main}.
Choisissez Full~Screen pour un écran de barre dédié;
  Windowed or Centred for a navigation station. \\
Préréglage de confort&
  \ui{Settings~> Comfort}.
  Set Day for initial testing; change to Night for the night watch. \\
Applications de lancement&
  \ui{Settings~> Applications}.
  Arrange your chart plotter, instruments, and other apps on page~1. \\
clavier virtuel&
  \ui{Settings~> Main} (touchscreen only).
  Tick, then restart FairWindSK. \\
\bottomrule
\end{longtable}

% ── A.12  Configuration file reference ───────────────────────
\section{Champs de configuration clés}
\label{app:config}

Advanced users and system integrators can edit \code{fairwindsk.json}
directement.
Cette section couvre les champs les plus utiles.

\subsection{Connexion}

\begin{lstlisting}
{
  "connection": {
    "server": "http://your-signalk-host:3000"
  }
}
\end{lstlisting}

\subsection{Mode fenêtre}

\begin{lstlisting}
{
  "main": {
    "windowMode": "fullscreen",
    "language": "en",
    "virtualKeyboard": false
  }
}
\end{lstlisting}

Valid values for \code{windowMode}: \code{windowed}, \code{centered},
\code{maximized}, \code{fullscreen}.
Valid values for \code{language}: \code{system}, \code{en}, \code{fr},
\code{es}, \code{it}.

\subsection{Unités}

\begin{lstlisting}
{
  "units": {
    "vesselSpeed": "kn",
    "windSpeed":   "kn",
    "distance":    "nm",
    "depth":       "mt"
  }
}
\end{lstlisting}

\subsection{Le chemin du signal K remplace}

\begin{lstlisting}
{
  "signalk": {
    "sog": "navigation.speedOverGround",
    "dpt": "environment.depth.belowTransducer",
    "notifications.pob": "notifications.mob"
  }
}
\end{lstlisting}

\subsection{Diagnostic}

\begin{lstlisting}
{
  "diagnostics": {
    "logLevel":          "off",
    "persistentLogs":    true,
    "interactionHistory": true,
    "email": "your-email@example.com"
  }
}
\end{lstlisting}

Valid log levels: \code{off}, \code{critical}, \code{warning},
\code{info}, \code{debug}, \code{full}.

% ── A.13  Troubleshooting installation ───────────────────────
\section{Dépannage des problèmes d'installation}
\label{app:troubleshooting-install}

\subsection{FairWindSK ne démarre pas après installation}

\begin{enumerate}
  \item Run the binary from a terminal to see error output:
        \code{FairWindSK} (Linux/macOS) or
        \code{build\textbackslash FairWindSK.exe} (Windows).
  \item Check that Qt WebEngine libraries are available.
        On Linux, run \code{ldd build/FairWindSK | grep "not found"}.
  \item On Raspberry~Pi, check for the virtual keyboard CMake workaround
        (Section~\ref{app:rpi-station}).
\end{enumerate}

\subsection{Aucune application n'apparaît sur le lanceur}

\begin{enumerate}
  \item Verify the Signal~K server address in
        \ui{Settings~> Connection}.
  \item Confirm the server exposes
        \code{/signalk/v1/apps/list} or \code{/skServer/webapps}.
  \item Check the Signal~K App Store --- at least one web app must be
installé sur le serveur pour qu'il apparaisse dans FairWindSK.
  \item Enable debug logging: add \code{debug=true} to
        \code{fairwindsk.ini} and relaunch.
Les journaux incluent les tentatives de connexion et les détails de découverte de l'application.
\end{enumerate}

\subsection{La construction échoue avec le module Qt manquant}

\begin{enumerate}
  \item Verify the Qt version: \code{qmake6 --version} or
        \code{cmake --find-package -DNAME=Qt6 -DCOMPILER\_ID=GNU \ldots}
  \item Check that all required packages are installed
        (Sections~\ref{app:linux} and~\ref{app:rpi-station}).
  \item If \code{CMAKE\_PREFIX\_PATH} is needed, pass it explicitly:
        \code{cmake -S . -B build -DCMAKE\_PREFIX\_PATH=/path/to/Qt/6.x.x/...}
\end{enumerate}

\subsection{FairWindSK ne démarre pas automatiquement sur Raspberry Pi}

\begin{enumerate}
  \item Verify the autostart file was installed:
\begin{lstlisting}
ls /etc/xdg/autostart/fairwindsk-startup.desktop
\end{lstlisting}
  \item Verify the \code{fairwindsk-startup} script is executable
et en PATH:
\begin{lstlisting}
which fairwindsk-startup
\end{lstlisting}
  \item Run the startup helper manually from a terminal to check for errors:
\begin{lstlisting}
fairwindsk-startup
\end{lstlisting}
  \item Ensure the desktop session supports XDG autostart
(le bureau Raspberry~Pi~OS par défaut fait).
\end{enumerate}
